\begin{abstract}
\noindent
Este trabajo aborda la planificación de una ruta turística de un día en Gran Canaria como un problema de orientación biobjetivo: elegir qué puntos de interés visitar, y en qué orden, para maximizar la puntuación turística y minimizar el tiempo total de la jornada (conducción más visitas), saliendo y volviendo al hotel en un máximo de 8 horas. El caso de estudio tiene 30 puntos de interés y dos hoteles, uno en el norte y otro en el sur.

Los tiempos de conducción se obtienen de la red viaria real de OpenStreetMap con OSMnx, con una penalización por la sinuosidad de las carreteras y un factor global calibrado con 12 trayectos medidos en Google Maps (error medio de 2,9 min). El problema se formula como un programa lineal entero mixto en variables de arco y se resuelve con una implementación propia de NSGA-II sobre rutas de longitud variable, con las reglas de dominancia de Deb para la restricción de tiempo y un operador de reparación que admite preferencias del turista (categorías, lugares descartados y una visita obligatoria).

NSGA-II se compara, con el mismo presupuesto de 20.100 evaluaciones y 15 semillas, con una heurística voraz, una búsqueda aleatoria y dos métodos de escalarización (suma ponderada y $\varepsilon$-restricción). En los dos hoteles obtiene el mayor hipervolumen en todas las semillas, entre un 10 \% y un 85 \% más que cada línea base, con diferencias significativas según el test de Wilcoxon ($p < 0{,}001$). La implementación se valida en los problemas ZDT1, ZDT2 y ZDT3. Los estudios de hiperparámetros y de sensibilidad muestran que la configuración elegida es razonable y que las rutas dependen sobre todo de la puntuación por minuto de cada lugar. De cada frente se extraen tres rutas recomendadas para ayudar a decidir.
\end{abstract}
